\documentclass{article}

% ==========================================
% 1. 核心中文支持 (必须)
% ==========================================
\usepackage[UTF8]{ctex} 

% ==========================================
% 2. 代码高亮设置
% ==========================================
\usepackage{listings}
\usepackage{xcolor}
\usepackage{geometry}
\geometry{a4paper, margin=1in}

\lstset{
    language=Python,               % 语言
    basicstyle=\ttfamily\small,    % 字体设置
    keywordstyle=\color{blue}\bfseries, % 关键字
    commentstyle=\color{green!50!black}, % 注释 (中文会在这里显示)
    stringstyle=\color{red},       % 字符串
    numbers=left,                  % 行号在左侧
    numberstyle=\tiny,
    stepnumber=1,
    frame=single,                  % 单线框
    breaklines=true,               % 自动换行
    % ----------------------------------------------------
    % 解决中文乱码的关键设置：
    extendedchars=false,           % 避免跨页时中文乱码
    keepspaces=true,               % 保持空格
    inputencoding=utf8,            % 明确指定输入编码
    % ----------------------------------------------------
    captionpos=b                   % 标题在底部
}

\begin{document}

\section{正文}
参见附录中的代码。

% ==========================================
% 附录
% ==========================================
\newpage
\appendix
\section{核心代码}

% -----------------------------------------------------------
% 注意事项：
% 1. caption={...} 里的文字如果有下划线 _，必须写成 \_
% 2. 后面 {} 里的文件路径，建议保持原样 (listings 会自动处理)
% -----------------------------------------------------------

% 1. 主程序
\subsection{运行脚本}
\lstinputlisting[language=Python, caption={主程序源码 (run.py)}]{codes/run.py}

% 2. 求解器
\subsection{优化求解器}
\lstinputlisting[language=Python, caption={求解器源码 (solver.py)}]{codes/solver.py}

% 3. 绘图工具 (注意 caption 里的下划线转义)
\subsection{绘图工具}
\lstinputlisting[language=Python, caption={绘图代码 (plots\_util.py)}]{codes/plots_util.py}

\end{document}